% !TEX root = ../ms-thesis.tex

\begin{abstract}
佛山市南海区具有较为丰厚的历史文化资源和旅游开发基础，但在实际的旅游景区建设中，文献记载、遗产名录中的文化记忆与游客能够接触到的旅游产品之间仍存在一定错位。识别这种错位，并判断哪些文化资源更具转化为旅游吸引物的条件，是区县尺度文旅融合研究需要回应的问题。基于此，本文以南海区为研究对象，从多源数据和空间分析角度识别旅游景区文旅融合的潜力关系。

研究建立''文化---旅游''双谱系对照框架。文化侧以地方典籍和相关文献为基础，通过文本整理、实体关系抽取和知识图谱构建，刻画南海区的文化记忆结构；旅游侧以多源 POI、平台评论数量、平台评分和景区等级等数据为基础，分析现有旅游供给与热度情况。两者之间，以官方认定且具有明确空间位置的物质文化载体作为主要耦合桥梁，并将非遗作为动态补充，构建文化记忆指数、官方认证指数、旅游热度指数和文化---旅游错位指数，对不同载体、片区和镇街的潜力状态进行比较。

研究发现，南海区旅游景区的文旅融合并不是简单的资源数量问题，而是文化资源、空间载体和旅游产品之间是否形成有效衔接的问题。一些文化记忆较深、官方认证较强的地区仍未充分转化为旅游热度，表现出较明显的''沉睡潜力''；部分旅游热度较高的现代景点，则存在本土文化叙事不足的情况。权重敏感性检验显示，主要分区结果在合理替代权重下具有一定稳定性；评论知识图谱补充检验也表明，旅游热度较高并不必然意味着游客已经充分感知地方文化。因此，文化基础和旅游热度同时较高的载体，更适合被理解为文旅融合潜力较大的优先深化对象，而不能简单判断为融合已经成熟。研究据此提出分层级、分片区的文旅融合提升建议，为南海区及同类区县的文化资源活化和旅游景区优化提供参考。

\thusetup{
  keywords = {文旅融合, 错位指数, 物质文化遗产, 知识图谱, 南海区},
}
\end{abstract}

\begin{abstract*}
Nanhai District in Foshan City has rich historical and cultural resources and a solid foundation for tourism development. However, in the construction of tourist attractions, a mismatch remains between the cultural memory recorded in documents and heritage inventories and the tourism products accessible to visitors. Identifying this mismatch and judging which cultural resources are more likely to be transformed into tourist attractions is a key issue for culture--tourism integration research at the district or county scale. This paper therefore takes Nanhai District as the research object and analyzes the potential of culture--tourism integration in tourist attractions from a multi-source data and spatial perspective.

This paper establishes a dual-spectrum comparative framework of ``culture'' and ``tourism.'' On the culture side, based on local chronicles and relevant literature, the structure of cultural memory in Nanhai District is characterized through text collation, entity-relationship extraction, and knowledge graph construction. On the tourism side, existing tourism supply and popularity are analyzed using data from multi-source points of interest (POI), online review counts, platform ratings, and attraction grades. Between the two, officially recognized tangible cultural carriers with clear spatial locations serve as the main coupling bridge, while intangible cultural heritage is used as a dynamic supplement. Based on these, the study further constructs a cultural memory index, an official recognition index, a tourism popularity index, and a culture--tourism mismatch index, thereby comparing potential states across different carriers, zones, and sub-districts or towns.

The findings reveal that culture--tourism integration in Nanhai's tourist attractions is not simply a matter of resource quantity, but rather whether an effective connection has been formed among cultural resources, spatial carriers, and tourism products. Some areas with deep cultural memory and strong official recognition have not yet been fully transformed into tourism popularity, exhibiting a pronounced ``dormant potential,'' while some modern attractions with high tourism popularity still lack sufficient local cultural narratives. Sensitivity tests suggest that the main zoning results remain relatively stable under reasonable alternative weights, and a supplementary review knowledge graph shows that high tourism popularity does not necessarily mean sufficient cultural recognition among visitors. Therefore, carriers with both strong cultural foundations and high tourism popularity should be understood as priority objects with high potential for further integration, rather than as evidence of mature integration. Based on these findings, this paper puts forward tiered and area-specific recommendations for enhancing culture--tourism integration.

\thusetup{
  keywords* = {culture–tourism integration, culture–tourism mismatch index, tangible cultural heritage, knowledge graph, Nanhai District},
}
\end{abstract*}
